nommen haben. Bis dahin hätte die Tafel aber mehr als 21000 Einträge. Noch wichtiger ist, dass man die dritten Differenzen keineswegs überall als verschwindend betrachten darf, wie \textsc{Leonelli} meint, der es für ausreichend hält, nur die zweiten zu berücksichtigen. Zu diesem Irrtum scheint ihn die von ihm berechnete Probe vom Anfang der Tafel verleitet zu haben, wo die dritten Differenzen tatsächlich verschwinden. Eine einfache Rechnung zeigt jedoch, dass diese dritten Differenzen im weiteren Verlauf \emph{zunehmen} und durchaus erheblich werden können, wenn man 14 Dezimalstellen angeben und keine kleineren Differenzen bei $P$ als 0,0005 verwenden will. Für $x = 2+\BareRoot{3}$, oder für $P = 0,5719$, wo die dritten Differenzen ihr Maximum erreichen, finden wir unter den genannten Voraussetzungen ihren Wert $= 6377$ Einheiten in der vierzehnten Dezimalstelle. \textsc{Leonelli} käme also bei Weitem nicht mit weniger als 5000 Einträgen aus. Wenn die zweiten Differenzen überall genügen sollen, müsste die Tafel so umfangreich sein, dass ihre Berechnung die Mühe keineswegs lohnen würde. Aber wozu überhaupt vierzehn Dezimalstellen? Rechnungen, die eine solche Genauigkeit erfordern, kommen nur äußerst selten vor. Für einen so seltenen Fall eine doch nicht sehr erhebliche Arbeitsersparnis zu gewinnen, ist von geringem Wert. Bei alltäglichen Rechnungen hingegen, für die sieben Dezimalstellen völlig ausreichen, wäre selbst eine mäßige, aber häufig wiederkehrende Erleichterung der Arbeit durchaus wertvoll. Um auch die nötige Bequemlichkeit zu bieten, müsste die Tafel dann so angelegt sein, dass überall auch die zweiten Differenzen verschwinden. Mit etwas mehr als 15000 Einträgen ließe sich das leicht erreichen. Sie ergäben nur ein schmales Bändchen. Für eine weitere Ausführung ist hier jedoch nicht der Ort.

Ein weiteres Beispiel dafür, wie oberflächlich \textsc{Leonelli} seinen Gegenstand behandelt hat, ist die auf S.\,60 angegebene Formel. Sie soll für einen vorgegebenen Wert von $R$, der zwischen zwei Werten der Tafel liegt, den entsprechenden Wert von $Q$ unter Berücksichtigung der zweiten Differenzen liefern. Diese Formel, deren Herleitung den Übersetzer nach eigenem Eingeständnis so viel vergebliche Mühe gekostet hat, ist ganz falsch. Statt
\[
2\delta(r'n - n) \quad \text{sollte nämlich stehen} \quad \frac{(D+d)\delta rn}{d}.
\]
